% Part: incompleteness
% Chapter: incompleteness-provability
% Section: rosser-thm

\documentclass[../../../include/open-logic-section]{subfiles}

\begin{document}

\olfileid[id]{inc}{inp}{ros}

\olsection{Teorema Rosser}

Dapatkah kita memodifikasi pembuktian G\"odel untuk memperoleh hasil yang
lebih kuat, dengan mengganti ``$\omega$-konsisten'' cukup dengan
``konsisten''? Jawabannya adalah ``ya,'' dengan menggunakan suatu siasat yang
ditemukan Rosser. Siasat Rosser ialah menggunakan predikat !!{derivability}
``termodifikasi'' $\ORProv_T(y)$ sebagai pengganti $\OProv[\Th{T}](y)$.

\begin{thm}
\ollabel{thm:rosser}
Misalkan $\Th{T}$ sebarang teori konsisten dan !!{axiomatizable} yang
memperluas $\Th{Q}$. Maka $\Th{T}$ tidak lengkap.
\end{thm}

\begin{proof}
Ingat bahwa $\OProv[\Th{T}](y)$ didefinisikan sebagai
$\lexists[x][\OPrf[\Th{T}](x,y)]$, dengan $\OPrf[\Th{T}](x,y)$
merepresentasikan relasi dapat diputuskan yang berlaku jika dan hanya jika
$x$ adalah bilangan G\"odel dari !!a{derivation} atas !!{sentence} bernomor
G\"odel~$y$. Relasi yang berlaku antara $x$ dan~$y$ jika $x$ adalah bilangan
G\"odel dari suatu \emph{refutasi} atas kalimat bernomor G\"odel~$y$ juga
dapat diputuskan. Misalkan $\fn{not}(x)$ adalah fungsi rekursif primitif yang
bekerja sebagai berikut: jika $x$ adalah kode suatu formula $!A$,
$\fn{not}(x)$ adalah kode~$\lnot !A$. Maka $\Refut[\Th{T}](x,y)$ berlaku jika
dan hanya jika $\Prf[\Th{T}](x,\fn{not}(y))$. Misalkan
$\ORefut[\Th{T}](x,y)$ merepresentasikannya. Jadi, jika
$\Th{T} \Proves \lnot !A$ dan $\delta$ adalah !!a{derivation} yang
bersesuaian, maka $\Th{Q} \Proves \ORefut[\Th{T}](\gn{\delta},\gn{!A})$.
Kita mendefinisikan $\ORProv[\Th{T}](y)$ sebagai
\[
\lexists[x][(\OPrf[\Th{T}](x,y) \land \lforall[z][(z < x \lif \lnot
  \ORefut[\Th{T}](z,y))])].
\]
Secara kasar, $\ORProv[\Th{T}](y)$ mengatakan ``terdapat pembuktian atas
$y$ dalam $\Th{T}$, dan tidak terdapat refutasi yang lebih pendek
atas~$y$.'' Dengan mengasumsikan $\Th{T}$ konsisten,
$\ORProv[\Th{T}](y)$ benar atas bilangan-bilangan yang sama dengan
$\OProv[\Th{T}](y)$; tetapi dari sudut pandang \emph{keterbuktian}
dalam~$\Th{T}$ (dan sekarang kita tahu bahwa kebenaran dan keterbuktian itu
berbeda!) keduanya memiliki sifat yang berbeda. Jika $\Th{T}$
\emph{tak} konsisten, keduanya \emph{tidak} berlaku atas bilangan-bilangan
yang sama!{} ($\ORProv[\Th{T}](y)$ sering dibaca sebagai ``$y$ terbukti
menurut Rosser.'' Karena, seperti baru dibahas, keterbuktian menurut Rosser
bukanlah suatu jenis keterbuktian khusus---dalam teori tak konsisten terdapat
!!{sentence}s yang terbukti tetapi tidak terbukti menurut Rosser---ungkapan
ini dapat membingungkan. Untuk menghindari kebingungan, Anda dapat membacanya
sebagai ``$y$ dapat di-shmove.'')

Berdasarkan lema titik tetap, terdapat suatu formula $!R_\Th{T}$ sedemikian
sehingga
\begin{equation}
  \Th{Q} \Proves !R_\Th{T} \liff \lnot \ORProv[\Th{T}](\gn{!R_\Th{T}}).
  \ollabel{RT}
\end{equation}
Berbeda dengan pembuktian \olref[1in]{thm:first-incompleteness}, di sini kita
menyatakan bahwa jika $\Th{T}$ konsisten, $\Th{T}$ tidak !!{derive}
$!R_\Th{T}$, dan $\Th{T}$ juga tidak !!{derive} $\lnot !R_\Th{T}$. (Dengan
kata lain, kita tidak memerlukan asumsi $\omega$-konsistensi.)

Pertama, mari kita tunjukkan bahwa $\Th{T} \Proves/ !R_{\Th{T}}$. Andaikan
sebaliknya, sehingga terdapat !!a{derivation} atas~$!R_{\Th{T}}$
dari~$T$; misalkan $n$ adalah bilangan G\"odel-nya. Maka
$\Th{Q} \Proves \OPrf[\Th{T}](\num{n},\gn{!R_{\Th{T}}})$, karena
$\OPrf[\Th{T}]$ merepresentasikan $\Prf[\Th{T}]$ dalam~$\Th{Q}$. Selain
itu, untuk setiap $k<n$, $k$ bukanlah bilangan G\"odel dari
!!a{derivation} atas $\lnot !R_{\Th{T}}$, karena $\Th{T}$ konsisten. Jadi,
untuk setiap $k<n$,
$\Th{Q} \Proves \lnot \ORefut[\Th{T}](\num{k},\gn{!R_{\Th{T}}})$.
Berdasarkan \olref[req][min]{lem:less-nsucc},
$\Th{Q} \Proves \lforall[z][(z < \num{n} \lif \lnot
\ORefut[\Th{T}](z,\gn{!R_{\Th{T}}}))]$. Dengan demikian,
\[
\Th{Q} \Proves \lexists[x][(\OPrf[\Th{T}](x,\gn{!R_{\Th{T}}}) \land
  \lforall[z][(z < x \lif \lnot \ORefut[\Th{T}](z,
  \gn{!R_{\Th{T}}}))])],
\]
tetapi itu tidak lain adalah $\ORProv[\Th{T}](\gn{!R_{\Th{T}}})$.
Berdasarkan \olref{RT}, $\Th{Q} \Proves \lnot !R_{\Th{T}}$. Karena
$\Th{T}$ memperluas $\Th{Q}$, juga
$\Th{T} \Proves \lnot !R_{\Th{T}}$. Kita telah mengasumsikan bahwa
$\Th{T} \Proves !R_{\Th{T}}$, sehingga $\Th{T}$ akan tak konsisten,
bertentangan dengan asumsi teorema.

Sekarang, mari kita tunjukkan bahwa
$\Th{T} \Proves/ \lnot !R_{\Th{T}}$. Sekali lagi, andaikan sebaliknya,
dan andaikan $n$ adalah bilangan G\"odel dari !!a{derivation}
atas~$\lnot !R_{\Th{T}}$. Maka $\Refut[\Th{T}](n,\Gn{!R_{\Th{T}}})$
berlaku, dan karena $\ORefut[\Th{T}]$ merepresentasikan
$\Refut[\Th{T}]$ dalam $\Th{Q}$,
$\Th{Q} \Proves \ORefut[\Th{T}](\num{n},\gn{!R_{\Th{T}}})$. Kita kembali
akan menunjukkan bahwa $\Th{T}$ akan tak konsisten karena teori itu juga
akan !!{derive}~$!R_{\Th{T}}$. Karena
\begin{align*}
\Th{Q} & \Proves !R_{\Th{T}} \liff \lnot \ORProv[\Th{T}](\gn{!R_{\Th{T}}}),
\intertext{dan karena $\Th{T}$ memperluas~$\Th{Q}$, cukup ditunjukkan bahwa}
\Th{Q} & \Proves \lnot \ORProv[\Th{T}](\gn{!R_{\Th{T}}}).
\end{align*}
!!^{sentence} $\lnot \ORProv[\Th{T}](\gn{!R_{\Th{T}}})$, yaitu
\begin{align*}
  \lnot & \lexists[x][(\OPrf[\Th{T}](x,\gn{!R_{\Th{T}}}) \land
    \lforall[z][(z < x \lif \lnot \ORefut[\Th{T}](z,
    \gn{!R_{\Th{T}}}))])],
  \intertext{secara logis ekuivalen dengan}
  & \lforall[x][(\OPrf[\Th{T}](x,\gn{!R_{\Th{T}}}) \lif
    \lexists[z][(z < x \land \ORefut[\Th{T}](z,
    \gn{!R_{\Th{T}}}))])].
\end{align*}
Kita menalar secara informal dengan menggunakan logika dan memanfaatkan
fakta-fakta mengenai hal-hal yang $\Th{Q}$ !!{derive}s. Andaikan $x$
sebarang dan $\OPrf[\Th{T}](x,\gn{!R_{\Th{T}}})$. Kita telah mengetahui
bahwa $\Th{T} \Proves/ !R_{\Th{T}}$, sehingga untuk setiap $k$,
$\Th{Q} \Proves \lnot \OPrf[\Th{T}](\num{k},\gn{!R_{\Th{T}}})$. Jadi,
untuk setiap $k$ diperoleh $\eq/[x][\num{k}]$. Khususnya, kita memiliki
(a)~$\eq/[x][\num{n}]$. Kita juga memiliki
$\lnot(\eq[x][\num{0}] \lor \eq[x][\num{1}] \lor \dots \lor
\eq[x][\num{n-1}])$, sehingga berdasarkan
\olref[req][min]{lem:less-nsucc}, (b)~$\lnot(x < \num{n})$. Berdasarkan
\olref[req][min]{lem:trichotomy}, $\num{n}<x$. Karena
$\Th{Q} \Proves \ORefut[\Th{T}](\num{n},\gn{!R_{\Th{T}}})$, kita memiliki
$\num{n}<x \land \ORefut[\Th{T}](\num{n},\gn{!R_{\Th{T}}})$, dan dari
hal itu $\lexists[z][(z < x \land \ORefut[\Th{T}](z,
\gn{!R_{\Th{T}}}))]$. Karena $x$ sebarang, sebagaimana diperlukan kita
memperoleh
\[
\lforall[x][(\OPrf[\Th{T}](x,\gn{!R_{\Th{T}}}) \lif
  \lexists[z][(z < x \land \ORefut[\Th{T}](z,
  \gn{!R_{\Th{T}}}))])].
\]
\end{proof}

\begin{prob}
Dua himpunan bilangan asli, $A$ dan $B$, disebut
\emph{tak terpisahkan secara komputabel} jika tidak ada himpunan
!!{decidable}~$X$ sedemikian sehingga $A \subseteq X$ dan
$B \subseteq \Complement{X}$ ($\Complement{X}$ adalah komplemen, yaitu
$\Nat \setminus X$, dari~$X$). Misalkan $\Th{T}$ adalah perluasan konsisten dan
!!{axiomatizable} dari~$\Th{Q}$. Andaikan $A$ adalah himpunan bilangan
G\"odel dari !!{sentence}s yang terbukti dalam~$\Th{T}$ dan $B$ adalah
himpunan bilangan G\"odel dari kalimat-kalimat yang terbantahkan
dalam~$\Th{T}$. Buktikan bahwa $A$ dan~$B$ tak terpisahkan secara
komputabel.
\end{prob}

\end{document}
